\section{Fazit und Ausblick}
\label{sec:fazit}

Das abschließende Kapitel fasst die wesentlichen Ergebnisse der vorliegenden Arbeit zusammen, beantwortet die eingangs formulierte Fragestellung und zeigt Ansatzpunkte für eine Weiterentwicklung des vorgestellten Prototyps auf.

\subsection{Zusammenfassung}

Gegenstand dieser Arbeit war die Konzeption und prototypische Implementierung eines Bestellanforderungsprozesses mit integriertem Genehmigungsworkflow auf der \gls{btp}. Ausgehend von der Problemstellung, dass Beschaffungsprozesse in vielen Unternehmen noch durch manuelle, papierbasierte oder E-Mail-gestützte Abläufe gekennzeichnet sind, wurde untersucht, wie sich durch die Kombination von \gls{cap}, SAP Fiori Elements und \gls{spa} ein durchgängiger, cloudnativer Geschäftsprozess auf der \gls{btp} realisieren lässt.

Die in der Einleitung definierten vier Teilziele wurden im Rahmen der Arbeit vollständig adressiert. Zunächst wurden in einer Anforderungsanalyse elf funktionale und vier nicht-funktionale Anforderungen systematisch erhoben und priorisiert (Teilziel~1). Darauf aufbauend wurde eine Anwendungsarchitektur entworfen, die auf einer klaren Drei-Schichten-Architektur basiert (Teilziel~2). Die Datenhaltungsschicht wird durch ein \gls{cds}-basiertes Datenmodell realisiert, das Bestellanforderungen mit ihren zugehörigen Positionen und Statusinformationen abbildet. Die Geschäftslogikschicht wurde mithilfe des \gls{cap} in Node.js implementiert und stellt die Fachlogik über \gls{odata}-Dienste bereit. Die Präsentationsschicht wurde auf Basis von SAP Fiori Elements umgesetzt, wobei der annotationsgetriebene Ansatz des \gls{cap} eine effiziente Generierung der Benutzeroberfläche ermöglichte, ohne dass umfangreicher Frontend-Code geschrieben werden musste.

Im Rahmen der prototypischen Implementierung (Teilziel~3) wurde der gesamte End-to-End-Prozess realisiert: Nach dem Erstellen einer Bestellanforderung über die Fiori-Elements-Oberfläche -- einschließlich des Hochladens von Dokumentenanhängen -- wird der Genehmigungsprozess in \gls{spa} automatisiert ausgelöst. Dem Genehmiger werden im Genehmigungsformular die vollständigen Kontextdaten einschließlich der Positionsliste, der Anhänge als Download-Links und einem Deep Link zur Fiori-Anwendung bereitgestellt. Das Ergebnis der Genehmigungsentscheidung wird über einen \gls{api}-Callback an den \gls{cap}-Dienst zurückgeschrieben, sodass der Status der Bestellanforderung in Echtzeit aktualisiert wird. Alle Statusänderungen werden in einer Statushistorie protokolliert, die den vollständigen Lebenszyklus jeder Bestellanforderung dokumentiert. Darüber hinaus können bereits genehmigte oder abgelehnte Anforderungen erneut bearbeitet und eingereicht werden.

Die Anwendung wurde als \gls{mta} auf Cloud Foundry bereitgestellt und anhand von elf Testfällen evaluiert (Teilziel~4). Dabei konnten alle elf funktionalen Anforderungen vollständig nachgewiesen werden.

Die Implementierung hat gezeigt, dass der annotationsgetriebene Ansatz des \gls{cap} in Kombination mit SAP Fiori Elements eine erhebliche Beschleunigung der Oberflächenentwicklung ermöglicht. Durch die deklarative Definition von Annotationen im \gls{cds}-Modell konnten Listen- und Detailansichten sowie Formulare ohne manuellen \gls{ui5}-Code generiert werden. Darüber hinaus hat sich \gls{spa} als geeignetes Werkzeug für die Low-Code-Modellierung von Genehmigungsworkflows erwiesen. Die grafische Prozessmodellierung ermöglicht eine schnelle Umsetzung von Workflow-Szenarien, während die nahtlose \gls{api}-Integration den bidirektionalen Datenaustausch mit der \gls{cap}-Anwendung sicherstellt.

Zusammenfassend lässt sich die eingangs formulierte Forschungsfrage positiv beantworten: Es ist sowohl technisch machbar als auch effizient, durch die Kombination von \gls{cap}, SAP Fiori Elements und \gls{spa} einen durchgängigen, cloudnativen Bestellanforderungsprozess mit Genehmigungsworkflow auf der \gls{btp} zu realisieren. Die Plattform stellt die notwendigen Bausteine bereit, um cloudnative Unternehmensanwendungen mit durchgängiger Prozessautomatisierung zu entwickeln. Die klare Trennung der Verantwortlichkeiten zwischen den einzelnen Diensten fördert dabei eine wartbare und erweiterbare Architektur, die den Anforderungen moderner Geschäftsprozessdigitalisierung gerecht wird.

\subsection{Ausblick}

Der im Rahmen dieser Arbeit entwickelte Prototyp bildet eine solide Grundlage, auf der weiterführende Entwicklungen aufbauen können. Im Folgenden werden ausgewählte Erweiterungsmöglichkeiten skizziert, die den praktischen Nutzen und die Produktionsreife der Anwendung erhöhen würden.

Eine naheliegende Erweiterung besteht in der Integration mit einem SAP S/4HANA-System \citep{metzger2021sap}. Im aktuellen Prototyp endet der Prozess mit der Genehmigung der Bestellanforderung. In einem produktiven Szenario sollte nach erfolgter Genehmigung automatisiert eine Bestellung im S/4HANA-Backend erzeugt werden, um den Beschaffungsprozess vollständig abzubilden. Hierzu könnten die von SAP bereitgestellten \gls{api}s des SAP S/4HANA-Systems genutzt werden, die sich über den SAP Destination Service in die \gls{btp}-Anwendung einbinden lassen.

Ein weiterer wesentlicher Aspekt ist die Einführung eines rollenbasierten Berechtigungskonzepts. Im vorliegenden Prototyp wird nicht zwischen der Rolle des Anfordernden und der des Genehmigenden unterschieden. In einer produktiven Umgebung sollte über den \gls{xsuaa}-Dienst eine feingranulare Rollenzuweisung erfolgen, sodass Anfordernde ausschließlich eigene Bestellanforderungen erstellen und einsehen können, während Genehmigende nur auf die ihnen zugewiesenen Genehmigungsaufgaben Zugriff erhalten.

Darüber hinaus wäre die Implementierung eines mehrstufigen Genehmigungsprozesses in Abhängigkeit vom Bestellwert von praktischem Nutzen. So könnte beispielsweise festgelegt werden, dass Bestellanforderungen ab einem bestimmten Schwellenwert -- etwa ab 10.000~EUR -- eine zusätzliche Genehmigung durch eine übergeordnete Instanz erfordern. \gls{spa} bietet mit seinen Entscheidungsregeln und Verzweigungslogiken die notwendigen Modellierungsmöglichkeiten, um derartige Szenarien umzusetzen.

Zur Verbesserung der Benutzererfahrung wäre zudem die Integration von Benachrichtigungsfunktionen sinnvoll. E-Mail- oder Push-Benachrichtigungen bei ausstehenden Genehmigungsaufgaben würden sicherstellen, dass Genehmigende zeitnah auf offene Anforderungen aufmerksam gemacht werden und Verzögerungen im Prozessablauf reduziert werden. Die im Prototyp implementierte Statushistorie bildet bereits eine Grundlage für die Nachvollziehbarkeit von Genehmigungsentscheidungen. Für ein umfassendes Compliance-Reporting wäre jedoch eine Erweiterung um die Protokollierung sämtlicher Datenänderungen an Kopfdaten und Positionen empfehlenswert.

Hinsichtlich der Bereitstellung und des Benutzerzugangs bietet sich eine Migration zu SAP Build Work Zone an. Durch die Einbindung der Anwendung in ein zentrales Launchpad könnten Nutzende über eine einheitliche Oberfläche auf die Bestellanforderungsanwendung sowie auf weitere Geschäftsanwendungen zugreifen, was die Akzeptanz und die Integration in bestehende Arbeitsabläufe verbessern würde.

Schließlich sollte für eine produktive Nutzung die Qualitätssicherung durch automatisierte Tests gestärkt werden. Das \gls{cap} stellt hierfür ein eigenes Testframework bereit \citep{herger2020cap}, mit dem sowohl Unit-Tests für einzelne Geschäftslogikkomponenten als auch Integrationstests für die \gls{odata}-Schnittstellen umgesetzt werden können. Eine Optimierung der Anwendung für große Datenvolumina, etwa durch serverseitige Paginierung und gezielte Datenbankindizierung, wäre ebenfalls im Hinblick auf die Produktionsreife empfehlenswert.

Insgesamt zeigt die vorliegende Arbeit, dass die \gls{btp} ein leistungsfähiges Fundament für die Entwicklung integrierter Geschäftsprozessanwendungen bietet \citep{newman2021microservices}. Die im Ausblick skizzierten Erweiterungen verdeutlichen zugleich das Potenzial, den vorgestellten Prototyp schrittweise zu einer produktionsreifen Lösung auszubauen, die den Anforderungen realer Unternehmensumgebungen gerecht wird.
